\section{方法}\label{sec:method}
本文的设计区分三个证据地位不同的角色：\textbf{生产求解器}是有界的“候选序~$\times$~在线溢出规划器”组合，也是唯一进入默认流水线的部分；\textbf{代价感知修复}是检验代价语义能否改进拓扑序的离线个案研究；\textbf{固定序精确规划器}是提供下界与证书的 CP-SAT 研究 oracle。

\subsection{生产求解器}
\paragraph{依赖前沿序。}大量同时就绪的 ALLOC 可能使单输入操作流长期饿死一个多输入消费者，并延长其已就绪操作数的驻留。依赖前沿序优先释放和可执行操作；仅剩分配节点时，完成最早可解锁消费者的就绪前驱组（图~\ref{fig:frontier-zh} 与算法~\ref{alg:canonical-zh}）。新增前沿序仅使用 DAG 结构与节点编号；组合内另保留 capfit-id、P1 序与 id-raw 三个既有候选序，本文不将整个组合重新表述为代价感知或模式无关的调度器。完整平局规则见补充材料第 3.1 节。

\begin{figure}[t]\centering
\includegraphics[width=.98\linewidth]{\ksFigure{frontier_mechanism.pdf}}
\caption{前驱组完成机制的示意算例。左：入度贪心序穿插无关的单输入流，三操作数消费者 $X$ 迟至第 9 步执行，累计驻留为 18 操作数步；右：依赖前沿序先完成就绪前驱组，使 $X$ 在第 4 步执行，累计驻留降至 6 操作数步（$-67\%$）。}
\label{fig:frontier-zh}\end{figure}

\begin{algorithm}[t]
\caption{Canonical 可扩展求解器}\label{alg:canonical-zh}
\begin{algorithmic}[1]
\Require 实例 $I$，任务 $p\in\{1,2,3\}$
\If{$p=1$} \State \Return \Call{P1MemoryOrder}{$I$} \EndIf
\State $\mathcal O\gets\{\Call{Frontier}{I},\Call{CapFitID}{I}\}$
\State 将 $\Call{P1Order}{I}$ 与 $\Call{IDRaw}{I}$ 加入 $\mathcal O$
\State $\mathcal R\gets\{\textsc{DistCost},\textsc{Adaptive25}\}$
\State $\mathcal H\gets\{0,5,40\}$（$p=2$），否则 $\{0,5,40,80,120\}$
\State $A^\star\gets\bot$；$k^\star\gets(+\infty,+\infty,+\infty)$
\ForAll{$(S,r,H)\in\mathcal O\times\mathcal R\times\mathcal H$}
  \State $A\gets\Call{BestFitPlan}{I,S,r,H}$
  \If{$A$ 可行}
    \State $(\mathrm{Tr},n,T)\gets\Call{Evaluate}{A}$
    \State $k\gets(\mathrm{Tr},n,T)$（$p=2$），否则 $(T,\mathrm{Tr},n)$
    \If{$k<k^\star$} \State $(A^\star,k^\star)\gets(A,k)$ \EndIf
  \EndIf
\EndFor
\State \Return $A^\star$ \Comment{拓扑序、地址偏移与溢出记录}
\end{algorithmic}
\end{algorithm}

\paragraph{地址分配与驱逐策略。}每个基础序使用 best-fit 连续地址分配，并比较距离/代价评分与自适应碎片策略。两者均为启发式，最终由完整工件的真实目标选优吸收其局部失误；精确评分、阈值、搬移边界与平局规则见补充材料第 3.1 节及表 S2。

\paragraph{真实目标组合选择。}P2/P3 分别在四个序、两种策略和 3/5 个重载窗口上形成固定规模组合，并按 $(\mathrm{Tr},n,T)$ 或 $(T,\mathrm{Tr},n)$ 选择完整工件。不可行成员被跳过，任何活跃量或溢出面积代理均不参与最终选择；窗口语义和完整工件生成过程见补充材料第 3.1--3.3 节。

\subsection{代价感知修复个案}
为探查代价语义进入拓扑序后的潜在余量，我们离线提出依赖合法的节点或链移动，重新规划并仅接受 P2 键严格改善的候选。该研究在不同实例上使用不同搜索器和预算，FlashAttention 小预算探针无收益，Matmul 未运行，因此它不是生产阶段或统一算法。搜索预算、负探针和逐案结果见补充材料第 3.4 节与第 4.5 节。

\subsection{固定序加权驻留间隙模型}
固定拓扑序后，每个缓冲区的相邻强制事件形成驻留间隙。保持间隙驻留消耗容量，中断则支付一次类别加权流量。二元间隙模型在逐缓存、逐切面的容量约束下最小化中断代价；其解还必须经过连续地址打包、溢出依赖生成与标准评测验证（图~\ref{fig:gap-model-zh}）。完整变量定义、容量切面、连续地址反例与模型编码见补充材料第 2.2--2.3 节和第 3.4 节。

\begin{figure}[t]\centering
\includegraphics[width=.98\linewidth]{\ksFigure{gap_model.pdf}}
\caption{加权驻留间隙模型。强制事件将缓冲区生命周期划分为可选间隙：保持驻留占用容量，中断则支付按类别加权的计费 $\kappa_b s_b$（有备份 $1\times$、生成 $2\times$）。CP-SAT 松弛给出流量下界；独立的连续地址打包与标准评测验证使达到下界的目标升级为固定序流量证书。}
\label{fig:gap-model-zh}\end{figure}

\begin{proposition}[固定序流量下界与证书]\label{prop:gap-zh}
对固定拓扑序，上述松弛的最优值是该序下一切合法连续地址溢出计划流量的下界；若某最优间隙选择可完成连续地址打包，且输出工件通过评测器验证并且流量等于该下界，则该工件在基准模型下是该固定序的流量最优计划。
\end{proposition}

证明将任意合法物理计划映射为容量松弛的可行间隙选择；完整证明以及“容量可行但无法连续打包”的反例见补充材料第 2.2--2.3 节。证书只优化固定序流量，不优化同流量下的次数或时间，也不证明跨拓扑序全局最优。

\paragraph{复杂度与部署边界。}在固定候选数下，可扩展路径为多项式实现，地址规划含保守的缓冲区二次项；间隙选择与后备连续打包最坏为指数复杂度。因此 CP-SAT 目前仅作研究 oracle，未来集成必须提供阈值/超时和经验证的回退；详细复杂度见补充材料第 2.5 节。
